% flow_ve_vs_poro.tex — §5.4.6 two long-time stress-relaxation channels
% Viscoelastic (Prony/Burgers) vs poroelastic (Terzaghi consolidation), common input & goal.
% Usage: \resizebox{\linewidth}{!}{% flow_ve_vs_poro.tex — §5.4.6 two long-time stress-relaxation channels
% Viscoelastic (Prony/Burgers) vs poroelastic (Terzaghi consolidation), common input & goal.
% Usage: \resizebox{\linewidth}{!}{% flow_ve_vs_poro.tex — §5.4.6 two long-time stress-relaxation channels
% Viscoelastic (Prony/Burgers) vs poroelastic (Terzaghi consolidation), common input & goal.
% Usage: \resizebox{\linewidth}{!}{% flow_ve_vs_poro.tex — §5.4.6 two long-time stress-relaxation channels
% Viscoelastic (Prony/Burgers) vs poroelastic (Terzaghi consolidation), common input & goal.
% Usage: \resizebox{\linewidth}{!}{\input{flow_ve_vs_poro.tex}}
\begin{tikzpicture}[
  every node/.style = {font=\small},
  N/.style  = {draw, rounded corners=3pt, align=center, thick,
               text width=6.4cm, minimum height=0.8cm, inner sep=4pt},
  Nin/.style= {N, fill=cyan!12, text width=8.4cm},
  Nve/.style= {N, fill=blue!10},
  Npo/.style= {N, fill=teal!14},
  Nout/.style={N, fill=purple!10, text width=8.4cm},
  ar/.style = {-Stealth, thick},
  var/.style= {-Stealth, thick, blue!55!black},
  par/.style= {-Stealth, thick, teal!55!black},
  lbl/.style= {font=\scriptsize\itshape, fill=white, inner sep=1.5pt, align=center},
]

\node[Nin] (in) {%
  \textbf{Common input} — instantaneous elastic interface stress ($t\!\to\!0$ limit from the UMAT)\\[1pt]
  $\sigma_{xx}^{0}$ (growth-induced residual stress, undrained / unrelaxed)};

% Route A — viscoelastic
\node[Nve, below left=1.15cm and -1.0cm of in] (a1) {%
  \textbf{Route A — viscoelastic channel}\\ (Burgers / Prony series)};
\node[Nve, below=0.45cm of a1] (a2) {%
  stress relaxation by sliding of the\\ polymer-chain (EPS) network};
\node[Nve, below=0.45cm of a2] (a3) {%
  time scale $\tau_{\mathrm{VE}}$\;$\to$\;
  decays to $\sim\!10\%$ of $\sigma_{xx}^{0}$};

% Route B — poroelastic
\node[Npo, below right=1.15cm and -1.0cm of in] (b1) {%
  \textbf{Route B — poroelastic channel}\\ (Terzaghi 1-D consolidation)};
\node[Npo, below=0.45cm of b1] (b2) {%
  growth-pressure gradient drives pore-fluid\\ drainage from the bulk surface};
\node[Npo, below=0.45cm of b2] (b3) {%
  consolidation scale $\tau=H^{2}/c_{v}$\;$\to$\\
  effective-stress halving \& dissipation};

\node[Nout, below=1.15cm of in |- a3.south, yshift=-0.7cm] (out) {%
  \textbf{Long-time steady state} ($t\!\to\!\infty$) — pressure-limited field\\[1pt]
  {\scriptsize set by the ratio of the two channels (Eq.\,5.7 long-time completion; Figs.\,5.19–5.20)}};

\draw[ar] (in.south) -- ++(0,-0.35) coordinate (split);
\draw[var] (split) -| (a1.north);
\draw[par] (split) -| (b1.north);
\draw[var] (a1) -- (a2);
\draw[var] (a2) -- (a3);
\draw[par] (b1) -- (b2);
\draw[par] (b2) -- (b3);
\draw[var] (a3.south) |- (out.west);
\draw[par] (b3.south) |- (out.east);

\begin{scope}[on background layer]
  \node[draw=blue!40, fill=blue!3, rounded corners=7pt, inner sep=6pt, fit=(a1)(a2)(a3)] {};
  \node[draw=teal!45, fill=teal!3, rounded corners=7pt, inner sep=6pt, fit=(b1)(b2)(b3)] {};
\end{scope}

\end{tikzpicture}
}
\begin{tikzpicture}[
  every node/.style = {font=\small},
  N/.style  = {draw, rounded corners=3pt, align=center, thick,
               text width=6.4cm, minimum height=0.8cm, inner sep=4pt},
  Nin/.style= {N, fill=cyan!12, text width=8.4cm},
  Nve/.style= {N, fill=blue!10},
  Npo/.style= {N, fill=teal!14},
  Nout/.style={N, fill=purple!10, text width=8.4cm},
  ar/.style = {-Stealth, thick},
  var/.style= {-Stealth, thick, blue!55!black},
  par/.style= {-Stealth, thick, teal!55!black},
  lbl/.style= {font=\scriptsize\itshape, fill=white, inner sep=1.5pt, align=center},
]

\node[Nin] (in) {%
  \textbf{Common input} — instantaneous elastic interface stress ($t\!\to\!0$ limit from the UMAT)\\[1pt]
  $\sigma_{xx}^{0}$ (growth-induced residual stress, undrained / unrelaxed)};

% Route A — viscoelastic
\node[Nve, below left=1.15cm and -1.0cm of in] (a1) {%
  \textbf{Route A — viscoelastic channel}\\ (Burgers / Prony series)};
\node[Nve, below=0.45cm of a1] (a2) {%
  stress relaxation by sliding of the\\ polymer-chain (EPS) network};
\node[Nve, below=0.45cm of a2] (a3) {%
  time scale $\tau_{\mathrm{VE}}$\;$\to$\;
  decays to $\sim\!10\%$ of $\sigma_{xx}^{0}$};

% Route B — poroelastic
\node[Npo, below right=1.15cm and -1.0cm of in] (b1) {%
  \textbf{Route B — poroelastic channel}\\ (Terzaghi 1-D consolidation)};
\node[Npo, below=0.45cm of b1] (b2) {%
  growth-pressure gradient drives pore-fluid\\ drainage from the bulk surface};
\node[Npo, below=0.45cm of b2] (b3) {%
  consolidation scale $\tau=H^{2}/c_{v}$\;$\to$\\
  effective-stress halving \& dissipation};

\node[Nout, below=1.15cm of in |- a3.south, yshift=-0.7cm] (out) {%
  \textbf{Long-time steady state} ($t\!\to\!\infty$) — pressure-limited field\\[1pt]
  {\scriptsize set by the ratio of the two channels (Eq.\,5.7 long-time completion; Figs.\,5.19–5.20)}};

\draw[ar] (in.south) -- ++(0,-0.35) coordinate (split);
\draw[var] (split) -| (a1.north);
\draw[par] (split) -| (b1.north);
\draw[var] (a1) -- (a2);
\draw[var] (a2) -- (a3);
\draw[par] (b1) -- (b2);
\draw[par] (b2) -- (b3);
\draw[var] (a3.south) |- (out.west);
\draw[par] (b3.south) |- (out.east);

\begin{scope}[on background layer]
  \node[draw=blue!40, fill=blue!3, rounded corners=7pt, inner sep=6pt, fit=(a1)(a2)(a3)] {};
  \node[draw=teal!45, fill=teal!3, rounded corners=7pt, inner sep=6pt, fit=(b1)(b2)(b3)] {};
\end{scope}

\end{tikzpicture}
}
\begin{tikzpicture}[
  every node/.style = {font=\small},
  N/.style  = {draw, rounded corners=3pt, align=center, thick,
               text width=6.4cm, minimum height=0.8cm, inner sep=4pt},
  Nin/.style= {N, fill=cyan!12, text width=8.4cm},
  Nve/.style= {N, fill=blue!10},
  Npo/.style= {N, fill=teal!14},
  Nout/.style={N, fill=purple!10, text width=8.4cm},
  ar/.style = {-Stealth, thick},
  var/.style= {-Stealth, thick, blue!55!black},
  par/.style= {-Stealth, thick, teal!55!black},
  lbl/.style= {font=\scriptsize\itshape, fill=white, inner sep=1.5pt, align=center},
]

\node[Nin] (in) {%
  \textbf{Common input} — instantaneous elastic interface stress ($t\!\to\!0$ limit from the UMAT)\\[1pt]
  $\sigma_{xx}^{0}$ (growth-induced residual stress, undrained / unrelaxed)};

% Route A — viscoelastic
\node[Nve, below left=1.15cm and -1.0cm of in] (a1) {%
  \textbf{Route A — viscoelastic channel}\\ (Burgers / Prony series)};
\node[Nve, below=0.45cm of a1] (a2) {%
  stress relaxation by sliding of the\\ polymer-chain (EPS) network};
\node[Nve, below=0.45cm of a2] (a3) {%
  time scale $\tau_{\mathrm{VE}}$\;$\to$\;
  decays to $\sim\!10\%$ of $\sigma_{xx}^{0}$};

% Route B — poroelastic
\node[Npo, below right=1.15cm and -1.0cm of in] (b1) {%
  \textbf{Route B — poroelastic channel}\\ (Terzaghi 1-D consolidation)};
\node[Npo, below=0.45cm of b1] (b2) {%
  growth-pressure gradient drives pore-fluid\\ drainage from the bulk surface};
\node[Npo, below=0.45cm of b2] (b3) {%
  consolidation scale $\tau=H^{2}/c_{v}$\;$\to$\\
  effective-stress halving \& dissipation};

\node[Nout, below=1.15cm of in |- a3.south, yshift=-0.7cm] (out) {%
  \textbf{Long-time steady state} ($t\!\to\!\infty$) — pressure-limited field\\[1pt]
  {\scriptsize set by the ratio of the two channels (Eq.\,5.7 long-time completion; Figs.\,5.19–5.20)}};

\draw[ar] (in.south) -- ++(0,-0.35) coordinate (split);
\draw[var] (split) -| (a1.north);
\draw[par] (split) -| (b1.north);
\draw[var] (a1) -- (a2);
\draw[var] (a2) -- (a3);
\draw[par] (b1) -- (b2);
\draw[par] (b2) -- (b3);
\draw[var] (a3.south) |- (out.west);
\draw[par] (b3.south) |- (out.east);

\begin{scope}[on background layer]
  \node[draw=blue!40, fill=blue!3, rounded corners=7pt, inner sep=6pt, fit=(a1)(a2)(a3)] {};
  \node[draw=teal!45, fill=teal!3, rounded corners=7pt, inner sep=6pt, fit=(b1)(b2)(b3)] {};
\end{scope}

\end{tikzpicture}
}
\begin{tikzpicture}[
  every node/.style = {font=\small},
  N/.style  = {draw, rounded corners=3pt, align=center, thick,
               text width=6.4cm, minimum height=0.8cm, inner sep=4pt},
  Nin/.style= {N, fill=cyan!12, text width=8.4cm},
  Nve/.style= {N, fill=blue!10},
  Npo/.style= {N, fill=teal!14},
  Nout/.style={N, fill=purple!10, text width=8.4cm},
  ar/.style = {-Stealth, thick},
  var/.style= {-Stealth, thick, blue!55!black},
  par/.style= {-Stealth, thick, teal!55!black},
  lbl/.style= {font=\scriptsize\itshape, fill=white, inner sep=1.5pt, align=center},
]

\node[Nin] (in) {%
  \textbf{Common input} — instantaneous elastic interface stress ($t\!\to\!0$ limit from the UMAT)\\[1pt]
  $\sigma_{xx}^{0}$ (growth-induced residual stress, undrained / unrelaxed)};

% Route A — viscoelastic
\node[Nve, below left=1.15cm and -1.0cm of in] (a1) {%
  \textbf{Route A — viscoelastic channel}\\ (Burgers / Prony series)};
\node[Nve, below=0.45cm of a1] (a2) {%
  stress relaxation by sliding of the\\ polymer-chain (EPS) network};
\node[Nve, below=0.45cm of a2] (a3) {%
  time scale $\tau_{\mathrm{VE}}$\;$\to$\;
  decays to $\sim\!10\%$ of $\sigma_{xx}^{0}$};

% Route B — poroelastic
\node[Npo, below right=1.15cm and -1.0cm of in] (b1) {%
  \textbf{Route B — poroelastic channel}\\ (Terzaghi 1-D consolidation)};
\node[Npo, below=0.45cm of b1] (b2) {%
  growth-pressure gradient drives pore-fluid\\ drainage from the bulk surface};
\node[Npo, below=0.45cm of b2] (b3) {%
  consolidation scale $\tau=H^{2}/c_{v}$\;$\to$\\
  effective-stress halving \& dissipation};

\node[Nout, below=1.15cm of in |- a3.south, yshift=-0.7cm] (out) {%
  \textbf{Long-time steady state} ($t\!\to\!\infty$) — pressure-limited field\\[1pt]
  {\scriptsize set by the ratio of the two channels (Eq.\,5.7 long-time completion; Figs.\,5.19–5.20)}};

\draw[ar] (in.south) -- ++(0,-0.35) coordinate (split);
\draw[var] (split) -| (a1.north);
\draw[par] (split) -| (b1.north);
\draw[var] (a1) -- (a2);
\draw[var] (a2) -- (a3);
\draw[par] (b1) -- (b2);
\draw[par] (b2) -- (b3);
\draw[var] (a3.south) |- (out.west);
\draw[par] (b3.south) |- (out.east);

\begin{scope}[on background layer]
  \node[draw=blue!40, fill=blue!3, rounded corners=7pt, inner sep=6pt, fit=(a1)(a2)(a3)] {};
  \node[draw=teal!45, fill=teal!3, rounded corners=7pt, inner sep=6pt, fit=(b1)(b2)(b3)] {};
\end{scope}

\end{tikzpicture}
